\section{Introduction}
\label{sec:intro}

\subsection{Background and motivation}

A time-to-digital converter quantizes the interval between two signal edges.
In an all-digital PLL it replaces the analog phase detector and charge pump,
which made TDCs standard in digital RF synthesis; they also appear in LiDAR
and time-of-flight imaging \cite{staszewski2005, henzlerbook}. The basic flash TDC, a delay line sampled
along its taps, has a resolution floor of one gate delay. In
\SI{180}{\nano\meter} CMOS that floor is several tens of picoseconds under
realistic loading (our own loaded stage measures \SI{97}{\pico\second} at
fan-out-of-four sizing), above the \SI{20}{\pico\second} this project
requires, so a sub-gate-delay technique is mandatory.

The Vernier principle reaches below the gate delay by racing two delay lines:
the START edge propagates through stages of delay $\ttau{1}$ while the STOP
edge chases it through faster stages of delay $\ttau{2}$. The effective
resolution is the difference $\tlsb=\ttau{1}-\ttau{2}$, which can be made
arbitrarily small in principle \cite{dudek2000}. The classical one-dimensional
implementation pays for this with hardware and latency: one stage pair and one
arbiter per code, so range and resolution fight each other. Ring topologies
reuse a short line to recover range at the cost of control complexity
\cite{yu2010}; passive interpolation attacks the problem from the matching side
\cite{henzler2008}. The two-dimensional Vernier of Liscidini, Vercesi and
Castello \cite{liscidini2009, vercesi2010} instead compares \emph{every} START
tap against \emph{every} STOP tap, so a grid of $N_X+N_Y$ stages covers
$\mathcal{O}(N_X N_Y)$ distinct intervals. We chose this architecture: it
decouples range from line length and suits a course budget measured in
transistor width.

\subsection{Objectives and boundary conditions}

The EE4615 brief sets hard requirements and a design-quality trade-off
(Table~\ref{tab:spec}). Every gate must be hand-designed in TSMC
\SI{180}{\nano\meter} BCD; standard cells are not allowed. Verification is by
Spectre simulation of the schematic, no layout. The course testbench fixes the
interface: a \SI{40}{\pico\second} RESET pulse, START and STOP inputs each
driven through $3\times$ minimum inverters, and 31 thermometer outputs
$Q_{1..31}$, each loaded by the \texttt{TDC\_out} cell (two $2\times$ inverters
and \SI{1}{\femto\farad}) \cite{ee4615tb}. Robustness must be demonstrated in
the five process corners TT, SS, FF, SF and FS \cite{ee4615corner}.

\begin{table}[H]
\centering
\small
\caption{Project requirements and quality metrics.}
\label{tab:spec}
\begin{tabular}{l l}
\toprule
\textbf{Requirement} & \textbf{Value} \\
\midrule
Output & 5 bit, thermometer code $Q_{1..31}$ \\
Resolution (LSB) & better than \SI{20}{\pico\second} \\
Range & full 32-code linear range \\
Process / supply & TSMC \SI{180}{\nano\meter} BCD, \SI{1.8}{\volt} \\
Corners & TT, SS, FF, SF, FS, all must pass \\
Standard cells & none; all cells hand-designed \\
Verification & schematic + Spectre only \\
\midrule
\textbf{Quality metric} & \textbf{Definition} \\
\midrule
Linearity & DNL, INL per corner \\
Energy & average energy per conversion \\
Area & total transistor width $\sumw$ \\
Figure of merit & energy per conversion step \\
\bottomrule
\end{tabular}
\end{table}

\subsection{Contributions and report structure}

The delivered converter reaches a calibrated resolution of
\SIrange{15.0}{15.85}{\pico\second} in all five corners with no missing codes.
Beyond meeting the specification, the project contributes three results that we
believe generalize: a quantified demonstration that fan-out-of-four sizing
cannot reach the required stage delays in this technology, a replica-load
sizing method for Vernier grids in which each tap drives on the order of twenty
unit gates, and a measured one-dimensional versus two-dimensional comparison
built from identical cells that isolates the architectural benefit from the
circuit design.

Section~\ref{sec:arch} develops the architecture and the operating point.
Section~\ref{sec:circuits} documents every cell with its transistor sizes.
Section~\ref{sec:method} describes the design iterations and the verification
setup. Section~\ref{sec:results} reports the corner, calibration, energy,
metastability and comparison results, and Section~\ref{sec:conclusion}
concludes.
